\section{FFN/MoE vs. Attention: diferencia esencial de los roles de cómputo}

\subsection{Dos paradigmas de cómputo}

\begin{importantbox}{Diferencia esencial entre Attention y FFN}
\begin{itemize}
  \item \textbf{Attention}: describe la interacción (interaction) \textbf{entre} tokens, es causal (causal mask); la query actual solo puede ver las key pasadas
  \item \textbf{FFN/MoE}: es una proyección (projection) a \textbf{nivel de token}, no hay interacción entre tokens; cada token pasa de forma independiente por la transformación FFN
\end{itemize}
\end{importantbox}

Esta diferencia tiene una implicación de ingeniería importante: como MoE/FFN es una operación token-level, entonces:
\begin{itemize}
  \item El enrutamiento es \textbf{per-token}: cada token elige de forma independiente sus top-K Expert
  \item Se pueden agrupar los tokens por Expert (grouping): los tokens enrutados al mismo Expert se agregan y se calculan en lote
  \item No involucra dependencias entre tokens, por lo que es naturalmente adecuado para la paralelización
\end{itemize}

\subsection{Análisis de la proporción de parámetros}

\begin{table}[H]
\centering
\caption{Proporción de parámetros: GQA vs. MoE/FFN}
\begin{tabular}{lccc}
\toprule
\textbf{Modelo} & \textbf{Proporción de GQA} & \textbf{Proporción de FFN/MoE} & \textbf{Nota} \\
\midrule
Qwen3-30B-A3B (128 Expert totales) & 2.9\% & 94\% & Parámetros totales: 30.5B \\
Qwen3-30B-A3B (8 Expert activados) & 27\% & 54\% & Parámetros activados: 3.3B \\
Qwen3-32B (Dense) & 18\% & 76\% & Parámetros totales: 32B \\
\bottomrule
\end{tabular}
\end{table}

\begin{knowledgebox}{FFN/MoE siempre concentra la mayor parte de los parámetros}
Ya sea Dense o MoE, la parte de FFN siempre ocupa la gran mayoría de los parámetros. En los parámetros totales de MoE, los 128 Expert aportan el 94\% de los parámetros. Incluso si solo se consideran los parámetros activados (8 Expert), FFN sigue representando el 54\%.
\end{knowledgebox}

\subsection{Resumen de la sección}

Attention se encarga de la interacción de información entre tokens, mientras que FFN/MoE se encarga de la transformación de características dentro de cada token. Ambas funciones son complementarias, pero la distribución de parámetros es sumamente desigual: FFN/MoE concentra la mayor parte, sin lugar a dudas.
